\writeSectionFrame{\presentationTitle}{Fazit}
\begin{frame}{Fazit}
	\begin{itemize}
 		\item Trennung der Verantwortlichkeiten als übergeordnetes Ziel
 		\item Das Model bleibt in fast allen Varianten der stabilste Teil
 		\item Evolution durch neue Anwendungsfälle und Probleme
 		\item Es gibt nicht das eine und richtige MVC
 	\end{itemize}
\end{frame}

\begin{frame}{Quo vadis MVC?}
	\begin{itemize}
 		\item Das C wird immer wichtiger, das M bleibt, das V löst sich auf
 		\item MVC wird zu MV*
 		\item Deklarative UIs und Reactive Programming
 		\item Micro-Frontends und Komponenten-basierte Architekturen
 	\end{itemize}
\end{frame}
\note{
	Das Model als reine Daten- und Geschäftslogik bleibt unverändert zentral.
	View wird immer flexibler und tendiert dazu, direkt auf Zustandsänderungen 
	zu reagieren, ohne dass ein expliziter Controller sie anweisen muss. C gewinnt
	an Bedeutung für zentrale Logik und für die Orchestrierung.
	
	 Die ursprüngliche Dreiteilung M-V-C ist oft zu starr. Viele moderne 
	 Architekturen sind flexibler und werden als MV* (Model-View-Whatever) 
	 bezeichnet. Die zugrunde liegenden Prinzipien der Trennung bleiben jedoch 
	 erhalten, auch wenn die Namen und genauen Verantwortlichkeiten variieren.
	 
	 UI-Frameworks setzen stark auf deklarative UIs.
	 
	 Große Anwendungen werden in kleinere, unabhängige Komponenten zerlegt 
	 (Micro-Frontends). Jede dieser Komponenten kann ihr eigenes internes MV* Muster
	 haben. Dies fördert die Skalierbarkeit und Wiederverwendbarkeit auf einer noch 
	 granulareren Ebene.
}

\begin{frame}{Häufige Paradigmen in modernen UIs}
	\begin{itemize}
 		\item Daten fließen in eine Richtung 
 		\item Änderungen lösen Re-Rendering aus
 		\item Unveränderliche Zustände
 	\end{itemize}
\end{frame}

\begin{frame}{MVCs Creed}
	\begin{itemize}
 		\item Grundlegendes Credo von MVC bleibt gleich trotz Namensänderungen und unterschiedlichen Umsetzungen 
 		\item Trennung von Sorgen als grundlegendes Bedürfnis
 		\item Lebt in neuen Formen weiter, immer angepasst an die aktuellen Bedürfnisse
 	\end{itemize}
\end{frame}


